\chapter*{Resumen}
\addcontentsline{toc}{chapter}{Resumen}

La gestión de proyectos de \emph{software} tradicional suele apoyarse en representaciones tabulares y listas que generan una alta carga cognitiva y fatiga visual, desviando la atención de los perfiles técnicos hacia tareas burocráticas y de microgestión \cite{meyer2019today}. Para dar solución a esta problemática, este Trabajo de Fin de Grado presenta Leaftask, una plataforma de gestión de proyectos que sustituye las vistas convencionales por metáforas visuales inmersivas basadas en grafos y árboles interactivos. Esta representación visual prioriza la claridad estructural y permite identificar dependencias, esfuerzos y el estado del trabajo de un vistazo \cite{ghoniem2004comparison}.

Además de la innovación en la interfaz de usuario, Leaftask integra un sistema multi-agente impulsado por Inteligencia Artificial. A diferencia de los asistentes convencionales, estos agentes operan como miembros proactivos del equipo: son capaces de ejecutar acciones de forma autónoma mediante herramientas (\emph{Tools}), reaccionar a eventos en la aplicación, comunicarse de forma asíncrona mediante chats y suspender sus flujos de trabajo a la espera de la interacción humana. Esto permite a los managers delegar muchas tareas rutinarias y de coordinación relacionadas con la gestión de proyectos, liberando tiempo para centrarse en tareas de mayor valor.

A nivel de ingeniería, la plataforma se ha construido sobre una base técnica robusta y escalable. Se ha diseñado e implementado una arquitectura de Monolito Modular orientada a eventos, aplicando los principios de la Arquitectura Hexagonal \cite{cockburn2005hexagonal} y la segregación de responsabilidades. La sincronización de datos entre los distintos \emph{Bounded Contexts} \cite{evans2004ddd} se resuelve bajo un modelo de consistencia eventual mediante la implementación de los patrones \emph{Command Query Responsibility Segregation (CQRS)} y \emph{Transactional Outbox/Inbox}.

Los resultados obtenidos validan la viabilidad técnica de la propuesta, logrando un producto funcional con una alta mantenibilidad y escalabilidad. En definitiva, el sistema proporciona un rendimiento óptimo en la interfaz gráfica y establece un entorno colaborativo donde la Inteligencia Artificial asume eficientemente la coordinación técnica rutinaria.